\documentclass[a4paper,12pt]{article}
\usepackage{geometry}
\geometry{letterpaper, margin=1in}
\usepackage{graphicx}
\usepackage{hyperref}
\usepackage{listings}
\usepackage{color}
\usepackage{tikz}

\title{
    \textbf{Campus Navigation and Safety Management System} \\
    \large Backend Architecture and Implementation Report
}
\author{
    \textbf{Project Team:} \\
    Muhammad Hanzala (Reg 232201020) \\
    Haseeb Nawaz (Reg 232201011) \\
    Saad Ali (Reg 232201015) \\[1em]
    \textbf{Instructor:} Sir Hussain Gillani \\
    \textbf{Course:} Web Technology
}
\date{\today}

\begin{document}

\maketitle

\begin{abstract}
This report details the implementation of the backend server for the Campus Navigation and Safety Management System (CN\_SMS). It covers the Flask-based architecture, RESTful API design, database schema, AI integration for safety analytics, and the administrative dashboard.
\end{abstract}

\tableofcontents
\newpage

\section{Introduction}
The CN\_SMS backend acts as the central intelligence hub, connecting the Android client app used by students/guards with the administrative web dashboard. It manages data consistency, provides navigation graphs, handles room bookings, and processes incident reports using AI.

\section{System Architecture}
The system is built on a micro-kernel architecture using Flask Blueprints to separate concerns.
\begin{enumerate}
    \item \textbf{Presentation Layer}: REST APIs (JSON) and Jinja2 Web View.
    \item \textbf{Business Logic Layer}: Controllers for Auth, Map, Incidents, and Bookings.
    \item \textbf{Data Layer}: SQLAlchemy ORM wrapping SQLite/PostgreSQL.
\end{enumerate}

\section{Database Schema (ERD)}
\begin{itemize}
    \item \textbf{User}: Stores credentials and roles (Admin, Security, Staff, Student).
    \item \textbf{MapNode/Edge}: Stores the graph needed for A* navigation on the client.
    \item \textbf{Incident}: Stores reports including images and AI analysis results.
    \item \textbf{Booking}: Manages room reservations with conflict detection.
\end{itemize}

\section{API Summary}
The backend exposes the following key endpoints:
\begin{itemize}
    \item \texttt{POST /api/auth/login}: JWT Authentication.
    \item \texttt{GET /api/map}: Retrieves the full campus graph.
    \item \texttt{POST /api/incidents/analyze}: Triggers Hugging Face AI to detect threats in images.
    \item \texttt{POST /api/bookings}: Creates new bookings with time-conflict validation (HTTP 409).
\end{itemize}

\section{AI Integration}
We integrate a Vision Transformer model (YOLOS-tiny) via the Hugging Face Inference API. The backend:
\begin{enumerate}
    \item Receives an image upload.
    \item Sends it to the Inference Endpoint.
    \item Processes detection tags (e.g., "fire", "person") to determine Severity (HIGH/MED/LOW).
    \item Falls back to a deterministic Mock analyzer if the external API is unreachable.
\end{enumerate}

\section{Conclusion}
The CN\_SMS backend successfully provides a robust, secure, and intelligent foundation for campus safety and navigation. It meets all specified requirements including modularity, security, and AI capabilities.

\end{document}
